\section{Introduction}\label{sec:intro}

Work-related accidents and diseases kill nearly three million people every year, 395 million workers sustain non-fatal injuries \cite{ilo2023}, and the economic burden approaches 4\% of global GDP \cite{ilo2023,lloyds2024}; in the United States alone it reached \$176.5 billion in 2023 \cite{nsc2023}. In heavy industry, personal protective equipment (PPE) is often the last barrier between a hazard and a worker, but a PPE mandate is only as effective as its enforcement. On most industrial sites, enforcement still means a human operator watching a wall of CCTV monitors.

That arrangement fails for a well-documented reason: vigilance is a rapidly depleting resource. Laboratory studies show detection performance collapsing within 20--35 minutes of continuous monitoring \cite{parasuraman1984,sawin1995}; surveillance-literature estimates suggest operators miss up to 95\% of on-screen activity after 22 minutes \cite{velastin2006}; and trained operators watching live feeds in diamond-processing plants detected only about half of targeted risk behaviors while generating high false-alarm rates \cite{donald2015}, with attention lapses accumulating over 90-minute sessions \cite{sandia2014}. A single camera produces 525{,}600 minutes of footage per year and an accident occupies roughly one of them. Manual surveillance pays off only if the operator is fully attentive during that exact minute.

Automated PPE detection promises to remove this bottleneck, and YOLO-family detectors dominate the field \cite{ultralytics2023,ppeyolo2024,edgeyolo2024}. Yet published systems are built and evaluated on construction-centric, lab-clean benchmarks; to our knowledge, none has been evaluated on imagery from a real, operating chemical plant. Two gaps produce that vacuum. First, the training data: public benchmarks such as Roboflow PPE-Detection \cite{roboflowppe} and the 17-class SH17 \cite{ahmad2024sh17} revolve around helmets and high-visibility vests and under-represent both the classes chemical operations depend on (gas masks, goggles) and the conditions of a phosphate-processing site (dust, harsh lighting, uniforms that generic models mistake for PPE). Second, few systems adapt to an individual site at all, even though that domain shift is precisely where generic detectors break down. This paper's central experiment measures the breakdown directly.

This paper addresses these gaps with a CCTV-retrofit PPE compliance system developed for, and evaluated at, OCP Group in Morocco. Its core claim is domain-specific rather than architectural: a detector fused and adapted for industrial and chemical conditions (oblique CCTV viewing angles, gas masks, site-specific uniforms) outperforms a detector trained only on the standard construction benchmark when both face chemical-plant imagery. Our contributions are:
\begin{itemize}
  \item \textbf{Real-plant central experiment.} A controlled comparison between a generic detector trained solely on the public construction benchmark \cite{roboflowppe} and our industrial-adapted model, on held-out OCP chemical-plant imagery and an industrial proxy set. The comparison, not a new architecture, is the claim.
  \item \textbf{Composite domain-adapted dataset.} A fusion of Roboflow PPE-Detection \cite{roboflowppe}, SH17 \cite{ahmad2024sh17}, the Roboflow gas masks corpus \cite{roboflowgasmasks}, and a self-annotated OCP corpus with weighted site-specific sampling, the dataset whose effect the central experiment measures.
  \item \textbf{Retrofit deployment model.} The pipeline runs on existing camera and network infrastructure; alerts reach the channels safety staff already use.
\end{itemize}

Section~\ref{sec:related} reviews related work; Section~\ref{sec:methods} details the dataset, comparison protocol, and system; Section~\ref{sec:results} presents results; Sections~\ref{sec:discussion} and~\ref{sec:conclusion} discuss limitations and conclude.
